\usepackage[T2A]{fontenc}
\usepackage[utf8]{inputenc}
\usepackage[russian]{babel}

% --- Шрифты и математика (без конфликта \Bbbk) ---
\usepackage{tempora}
\usepackage{amsmath}
\let\Bbbk\relax  % Устраняет конфликт между amssymb и newtxmath
\usepackage{amssymb}
\usepackage{newtxmath}
\usepackage{microtype}

% --- Геометрия страницы и колонок под формат конференции ---
\usepackage[left=20mm, right=20mm, top=25mm, bottom=25mm]{geometry}
\setlength{\columnsep}{6mm}  % Оптимальный зазор между колонками

% --- Межстрочный интервал и абзацы ---
\usepackage{setspace}
\singlespacing               % Одинарный интервал для двухколоночной верстки
\usepackage{indentfirst}
\setlength{\parindent}{0.6cm} % Компактный отступ красной строки (не раздувает колонку)
\usepackage{hyperref}
\usepackage{orcidlink} % при наличии пакета в TeX Live, либо стандартный \href
% --- Графика и таблицы ---
\usepackage{graphicx}
\graphicspath{{fig/}}
\usepackage{booktabs}
\usepackage{array}
\usepackage{float}
\usepackage{multirow}

% --- Подписи к рисункам и таблицам ---
\usepackage{caption}
\DeclareCaptionLabelSeparator{endash}{\space---\space}
\captionsetup[figure]{labelsep=endash, justification=centering, font=footnotesize}
\captionsetup[table]{labelsep=endash, justification=raggedright, singlelinecheck=false, font=footnotesize}

% --- Заголовки разделов для двух колонок ---
\usepackage{titlesec}

\titleformat{\section}[hang]
    {\raggedright\bfseries\normalsize}
    {\thesection.}
    {0.4em}
    {}

\titleformat{\subsection}[hang]
    {\raggedright\itshape\normalsize}
    {\thesubsection.}
    {0.4em}
    {}

\titlespacing*{\section}{0pt}{8pt plus 2pt minus 1pt}{4pt plus 1pt minus 1pt}
\titlespacing*{\subsection}{0pt}{6pt plus 1pt minus 1pt}{3pt plus 1pt minus 1pt}

% --- Подсветка кода (компактная для узких колонок) ---
\usepackage{xcolor}
\usepackage{listings}
\definecolor{codegreen}{rgb}{0,0.6,0}
\definecolor{codegray}{rgb}{0.5,0.5,0.5}
\definecolor{codepurple}{rgb}{0.58,0,0.82}
\definecolor{backcolour}{rgb}{0.96,0.96,0.96}

\lstdefinestyle{gostcode}{
    backgroundcolor=\color{backcolour},
    commentstyle=\color{codegreen},
    keywordstyle=\color{magenta},
    numberstyle=\tiny\color{codegray},
    stringstyle=\color{codepurple},
    basicstyle=\ttfamily\scriptsize, % Уменьшенный шрифт, чтобы строки не ломались
    breaklines=true,
    captionpos=t,
    numbers=left,
    numbersep=4pt,
    tabsize=2,
    frame=single,
    keepspaces=true,
    extendedchars=true
}
\lstset{style=gostcode}

% --- Ссылки и URL (с корректным переносом в узких строках) ---
\usepackage{url}
\def\UrlBreaks{\do\/\do-\do_}
\usepackage{hyperref}
\hypersetup{
    colorlinks=true,
    linkcolor=black,
    filecolor=black,
    urlcolor=blue,
    citecolor=black
}
% --- Настройка колонтитулов ---
\usepackage{fancyhdr}
\pagestyle{fancy}
\fancyhf{} % Очистка текущих настроек

% Линейка под верхним колонтитулом
\renewcommand{\headrulewidth}{0.4pt}
\renewcommand{\footrulewidth}{0pt}

% Верхние колонтитулы для четных (E) и нечетных (O) страниц
\fancyhead[LE,RO]{\small\thepage}
\fancyhead[RE]{\small\itshape И.\,О.~Фамилия} % Ваши инициалы и фамилия
\fancyhead[LO]{\small\itshape Аппаратно-программный комплекс первичного поиска сигналов GPS L1}

% На первой странице колонтитулы отключаются (стиль plain)
\fancypagestyle{plain}{%
  \fancyhf{}
  \renewcommand{\headrulewidth}{0pt}
}
